% ---
% Inicia os anexos
% ---
\begin{anexosenv}
    \chapter{Código Exercício 4}\label{cap:ex4-codigo}
    
    \begin{lstlisting}[language=Python, breaklines=true, basicstyle=\tiny]
    import numpy as np
    import matplotlib.pyplot as plt
    import scipy.special
    from scipy.special import eval_jacobi
    from scipy.interpolate import BarycentricInterpolator

    def equidist_nodes_triangle(N):
        points = []

        for i in range(N + 1):
            for j in range(N + 1 - i):

                # Barycentric coordinates
                l1 = i / N if N > 0 else 0.0
                l2 = j / N if N > 0 else 0.0
                l3 = 1.0 - l1 - l2

                # Triangle vertices
                v1 = np.array([-1.0, -1.0])
                v2 = np.array([ 1.0, -1.0])
                v3 = np.array([-1.0,  1.0])

                # Cartesian coordinates
                p = l1 * v2 + l2 * v3 + l3 * v1
                points.append(p)

        return np.array(points)

    def get_gll_nodes(N):
        internal_nodes, _ = scipy.special.roots_jacobi(N - 1, 1, 1)
        r = np.concatenate([[-1.0], internal_nodes, [1.0]]) # elemento de referência
        return r

    def blend(a, b, c, alpha):
        num = (4*a*b)**alpha
        den = num + c**alpha
        
        if den == 0: 
            return 0
        
        return num/den

    def warped_nodes_triangle(N):
        r_list, s_list = [], []
        
        # configura a função warp
        alpha = 2
        uniform = np.linspace(-1, 1, N+1)
        gll_nodes = get_gll_nodes(N)
        diff_gll_uniform = gll_nodes - uniform
        warp = BarycentricInterpolator(uniform, diff_gll_uniform)
        
        for i in range(N+1):
            for j in range(N - i + 1):
                k=N-i-j
                lambda_1, lambda_2, lambda_3 = i/N, j/N, k/N
                r_eq = 2*lambda_2 - 1
                s_eq = 2*lambda_3 - 1
                
                zeta_1 = lambda_2 - lambda_1 
                zeta_2 = lambda_3 - lambda_2 
                zeta_3 = lambda_1 - lambda_3 
                w1 = warp(zeta_1)
                w2 = warp(zeta_2)
                w3 = warp(zeta_3)
                
                b1 = blend(lambda_1, lambda_2, lambda_3, alpha)
                b2 = blend(lambda_2, lambda_3, lambda_1, alpha)
                b3 = blend(lambda_3, lambda_1, lambda_2, alpha)

                dr = b1 * w1 + b2 * (-w2)
                ds = b2 * w2 + b3 * (-w3)
                
                r_list.append(r_eq + dr)
                s_list.append(s_eq + ds)

        return np.array(r_list), np.array(s_list)

    def simplex_2D(i, j, r, s):
        a = 2.0*(1.0 + r)/(1.0 - s) - 1.0
        b=s
        
        # singularidade em s=1 (vértice superior)
        mask = np.abs(s - 1.0) < 1e-12
        a[mask] = -1.0
        h1 = eval_jacobi(i, 0, 0, a)
        h2 = eval_jacobi(j, 2*i+1, 0, b)
        
        return np.sqrt(2.0) * h1 * h2 * (1.0 - b)**i

    def vandermonde2D(N, r, s):
        Np = (N+1)*(N+2)//2
        V = np.zeros((len(r), Np))
        col = 0
        for i in range(N+1):
            for j in range(N - i + 1):
                V[:, col] = simplex_2D(i, j, r, s)
                col += 1
        return V

    plt.rcParams.update({
        "font.size": 14
    })

    fig, axes = plt.subplots(1, 3, figsize=(16, 6))
    N_list = [3,5,8]

    for i, ax in enumerate(axes.flat):
        # r, s = warped_nodes_triangle(N_list[i])
        # ax.scatter(r, s, c='black')
        
        points = equidist_nodes_triangle(N_list[i])
        ax.scatter(points[:,0], points[:,1], c='black')
        
        tri = np.array([[-1,-1], [1,-1], [-1,1], [-1,-1]])
        ax.plot(tri[:,0], tri[:,1], '--', lw=2, color='gray' )
        ax.set_ylabel(r'$s$')
        ax.set_xlabel(r'$r$')
        ax.set_title(f'Nós Equidistantes para N={N_list[i]}')

    plt.tight_layout()
    plt.show()
    \end{lstlisting}

    \chapter{Código Exercício 6}\label{cap:ex6-codigo}
    
    \begin{lstlisting}[language=Python, breaklines=true, basicstyle=\tiny]
    import numpy as np
    import matplotlib.pyplot as plt

    # Malha
    N = 3
    x = np.linspace(-1.0, 1.0, N)
    y = np.linspace(-1.0, 1.0, N)
    dx = x[1] - x[0]
    dy = y[1] - y[0]
    X, Y = np.meshgrid(x, y)

    # Gradiente Az
    Az = X**2 + Y**2
    dAz_dy, dAz_dx = np.gradient(Az, dy, dx)

    # Campo magnético B = grad(Az) x z_hat
    Bx = dAz_dy
    By = -dAz_dx

    # Plot
    plt.figure(figsize=(6,6))
    plt.quiver(X, Y, Bx, By, scale_units='xy', scale=4,  angles='xy')

    plt.xlabel("x")
    plt.ylabel("y")
    plt.title(r"$\mathbf{B}=\nabla A_z\times\hat{z} \quad$  N = 3")
    plt.axis("equal")
    plt.grid(True)
    plt.ylim(-1.5,1.5)
    plt.xlim(-1.5,1.5)
    plt.tight_layout()

    plt.show()
    \end{lstlisting}

\end{anexosenv}
